% ECONOMICS_PROBLEM: {"id": "Q53-427", "title": "Causal pollution dose–response and avoidance", "jel_code": "Q53", "source_id": "OP-0118"}
% !TeX program = xelatex
\documentclass[11pt,a4paper]{article}
\usepackage[margin=23mm]{geometry}
\usepackage{fontspec}
\setmainfont{Latin Modern Roman}
\usepackage{amsmath,amssymb,mathtools}
\usepackage{xcolor,enumitem,booktabs,longtable,array,xurl,hyperref}
\definecolor{muted}{HTML}{53636C}
\hypersetup{colorlinks=true,urlcolor=blue,linkcolor=blue}
\setlength{\parindent}{0pt}
\setlength{\parskip}{5pt}
\newcommand{\R}{\mathbb R}
\newcommand{\E}{\mathbb E}
\newcommand{\N}{\mathbb N}
\newcommand{\ind}{\mathbf 1}
\newcommand{\Var}{\operatorname{Var}}
\newcommand{\Cov}{\operatorname{Cov}}
\newcommand{\supp}{\operatorname{supp}}
\DeclareMathOperator*{\argmin}{arg\,min}
\DeclareMathOperator*{\argmax}{arg\,max}
\newcommand{\entrymeta}[3]{{\sffamily\small\color{muted}\textbf{\texttt{#1}}\quad|\quad #2\par #3\par}}
\newcommand{\Problemref}[1]{\texttt{#1}}
\newcommand{\JELref}[2]{\texttt{#2} (JEL \texttt{#1})}
\begin{document}
\section*{Causal pollution dose–response and avoidance}
\textbf{Problem ID:} \texttt{Q53-427}\quad \textbf{JEL:} \texttt{Q53}\par
% BEGIN ECONOMICS STATEMENT
\label{problem:OP-0118}
\label{atlas:prob:C8}

\noindent\textbf{Original identifier:} C8 (atlas item 118). \textbf{Field:} Environmental and climate economics. \textbf{Classification:} editorial research family with established restricted solutions; current open status not certified.

\subsection*{Definitions and assumptions}

For individuals $i$ over a fixed horizon, let $\bar p$ be an ambient history of a vector of pollutants, with each concentration in declared units (for example $\mu\mathrm g/\mathrm m^3$). Avoidance regime $a_i(\bar p)$ includes location, filtration and time allocation; personal dose is $e_i(\bar p,a_i)$ in a specified physical unit. Health/productivity potential outcome $Y_i(\bar p,a)$ depends on exposure and avoidance. Define the controlled response $\mu(\bar p,a)=\mathbb E[Y_i(\bar p,a)]$ for a fixed population, and the policy response $\mu^{\rm pol}(\bar p)=\mathbb E[Y_i(\bar p,a_i(\bar p))]$. Consistency requires a fully specified intervention, including which pollutants and exposure components change.

\subsection*{Formal research question}

Characterize the identifiable dose-response surface and total welfare loss under endogenous avoidance, correlated pollutant mixtures and imperfect exposure measurement. For an intervention from $\bar p_0$ to $\bar p_1$, define
\[
\Delta W=\mathbb E\!\left[
v_i(Y_i(\bar p_1,a_i(\bar p_1)))-k_i(a_i(\bar p_1))
-v_i(Y_i(\bar p_0,a_i(\bar p_0)))+k_i(a_i(\bar p_0))
\right],
\]
with $v_i$ a declared monetary valuation and $k_i$ real avoidance costs in the same currency year. Specify a sampling class and quasi-experimental assignment mechanism, then derive sharp bounds and uniform inference for $\Delta W$ or selected components. Determine what dose ranges and mixture directions are supported by the instrument, and what additional smoothness, exclusion and transport restrictions are required outside that support.

\subsection*{Interpretation and scope}

Wind direction or inversions can shift pollution but may also affect temperature, activity and infectious disease; exclusion must be justified for the actual health target. Instrumental variation generally identifies local or instrument-weighted responses, not an arbitrary full nonlinear surface. Avoidance is a mediator: controlling for it changes a total-policy estimand into another contrast and may create selection bias. Ambient monitors measure locations rather than individual inhaled dose, and mobility creates additional selection. Acute mortality effects among elderly patients differ from chronic lifetime harms and productivity effects. Costs of defensive spending, medical use and mortality valuation need coherent welfare accounting to avoid double counting.

\subsection*{Source audit and status}

[S1] is completed as the verified recession-shock infant-mortality article; Chay--Greenstone also have a distinct 2003 Clean Air Act working paper, so input identity is not uniquely determined. [S2] studies California infant health; [S3] estimates acute PM2.5 effects among US elderly using wind. AirNow monitoring supports exposure data, not causal identification. The full mixture-and-avoidance welfare target is editorial; these established studies do not certify its exact 2026 open status.

\subsection*{References}

\begin{enumerate}[label={[\textnormal{S}\arabic*]},leftmargin=*]

\item Kenneth Y. Chay and Michael Greenstone (2003). \emph{The Impact of Air Pollution on Infant Mortality: Evidence from Geographic Variation in Pollution Shocks Induced by a Recession}. Quarterly Journal of Economics 118(3), 1121–1167. \url{https://doi.org/10.1162/00335530360698513}. \textbf{Locator:} Publisher or institutional abstract and bibliographic record. \textbf{Support and limits:} Recession-induced TSP variation and infant mortality; same author-year also has a Clean Air Act working paper, so title completion is disclosed.

\item Janet Currie and Matthew Neidell (2005). \emph{Air Pollution and Infant Health: What Can We Learn from California’s Recent Experience?}. Quarterly Journal of Economics 120(3), 1003–1030. \url{https://academic.oup.com/qje/article-abstract/120/3/1003/1841459}. \textbf{Locator:} Publisher or institutional abstract and bibliographic record. \textbf{Support and limits:} California infant-health pollution estimates; does not identify all pollutants, populations, exposure histories or avoidance costs.

\item Tatyana Deryugina, Garth Heutel, Nolan H. Miller, David Molitor and Julian Reif (2019). \emph{The Mortality and Medical Costs of Air Pollution: Evidence from Changes in Wind Direction}. American Economic Review 109(12), 4178–4219. \url{https://doi.org/10.1257/aer.20180279}. \textbf{Locator:} Publisher or institutional abstract and bibliographic record. \textbf{Support and limits:} Acute PM2.5 effects among US elderly using wind variation; requires exclusion and does not establish a universal chronic dose-response.

\end{enumerate}

\subsection*{Common conventions: Source mathematical conventions}

Unless an entry states otherwise, agent and firm sets are finite. State and action spaces are measurable, strategies and outcome maps are measurable, probability laws are specified, and expectations exist for the targets under discussion. Infinite-horizon discount factors lie in $(0,1)$ and discounted objectives are finite. Norms, tolerances and complexity input sizes must be specified for computational comparisons. Constants or primitives that appear locally are inputs, not empirical facts asserted by this catalogue.

\textbf{Identification.} A statistical model $m\in\mathcal M$ implies an observable law $P_m$ and target $\theta(m)$. At law $P$, its identified set is
\[
\Theta(P;\mathcal M)=\{\theta(m):m\in\mathcal M,\ P_m=P\}.
\]
Point identification holds when this set is a singleton, or equivalently when $P_{m_1}=P_{m_2}$ implies $\theta(m_1)=\theta(m_2)$ for all compatible models. Sharp bounds mean bounds attained, or approached when the boundary is not attained, by compatible models. Writing this set is a definition, not a solution: the substantive task is to establish informative restrictions and derive or estimate the set. An unrestricted-confounding model may give only trivial bounds.

\textbf{Inference.} Identification, estimability and valid uncertainty quantification are separate questions. Fix $\alpha\in(0,1)$. A confidence procedure $C_n$ over a declared law class $\mathcal P$ has asymptotic uniform coverage at level $1-\alpha$ when
\[
\liminf_{n\to\infty}\inf_{P\in\mathcal P}\Pr_P\{\theta(P)\in C_n\}\geq1-\alpha.
\]
For set-identified targets the coverage event must be defined separately, for example coverage of the full identified set. If an entry uses identification-indexed models instead of uniquely indexed laws, the same coverage convention is applied over those models. Pointwise limits do not establish uniform validity, and asymptotic validity does not guarantee finite-sample precision. Sample sizes, dependence structures and weak-identification sequences are part of each application.

\textbf{Counterfactuals.} $Y_i(a)$ is a potential outcome under a specified intervention, including its timing, implementation and versions. It is not $Y_i$ conditional on voluntarily choosing $a$. A full intervention vector permits interference; shorthand scalar treatment notation does not silently impose no interference. Assignment, selection, attrition, anticipation and measurement must be specified. A claim about an instrument's local effect is not automatically a population policy effect.

\textbf{Economic equilibrium.} A counterfactual model includes best responses, constraints, feasibility, market clearing, state transitions and expectations consistent with the stated information. When several equilibria exist, either a declared selector is an input or the target is set-valued. Comparisons across policy regimes must use the stated selection convention. Reduced-form relationships are not presumed invariant after an equilibrium-changing intervention.

\textbf{Optimization and welfare.} Optimization is over the stated feasible set. If attainment is not established, $\sup$ or $\inf$ and $\varepsilon$-optimal choices replace an asserted maximizer or minimizer. For example, with value $V=\sup_{a\in\mathcal A}W(a)<\infty$, an $\varepsilon$-optimal set is $\{a\in\mathcal A:W(a)\geq V-\varepsilon\}$ for $\varepsilon>0$. Benefits, costs, risk and health or environmental outcomes can be aggregated only using declared units and conversion weights. Interpersonal utility weights, the treatment of fiscal transfers and foreign profits, discounting, and the behavioral welfare criterion are explicit inputs. A comparative-static derivative additionally requires existence and appropriate differentiability of the selected counterfactual.

\textbf{Sources.} Local labels [S1], [S2], and so forth restart in every question. Locators identify the relevant abstract, model, section or result at the level actually checked; a bibliographic or abstract-level verification is not represented as inspection of an entire proof. Working papers are distinguished from later publications and changing versions. Source summaries are paraphrased. All URL checks and editorial formulations refer to the audit date stated above.
% END ECONOMICS STATEMENT
\end{document}
